%% PA sequence extension and specified primitive-recursive graph formulas.
%% This supplies an internal-arithmetic prerequisite for the proof predicate.

\begin{annexitem}{annex:p5.internal-pr-graphs}{thm:p05-pa-pr-graphs}{Primitive recursive graphs verified in PA}
We prove \cref{thm:p05-pa-pr-graphs} using the seven arithmetic axioms and
the induction scheme defining \(\PA\).  A statement made ``in \(\PA\)''
below means a derivation with all its displayed free parameters universally
closed.  Induction may therefore be applied to the formulas with those
parameters.  No claim about all codes is inferred from an external computation.
Polynomial juxtaposition abbreviates multiplication, natural constants stand
for their numeral terms, and intervals in the explanatory prose abbreviate
the corresponding arithmetic inequalities.  All auxiliary formulas use
fresh bound variables before their arguments are substituted.

\paragraph{Elementary arithmetic used in the construction.}
The arithmetic manipulations below take place in \(\PA\).  Here are the
induction arguments needed to fix their standing.  Induction gives
\(0+x=x\) and \(Sx+y=S(x+y)\); induction on the second addend then gives
commutativity, and induction on the last addend gives associativity.
Induction on \(z\), using successor injectivity in its step, gives
\(x+z=y+z\to x=y\).  The multiplication recursion, together with these
addition identities, gives distributivity by induction on the recursive
argument; it gives \(0\times x=0\) and
\((Sx)\times y=x\times y+y\) by induction on \(y\), and hence
commutativity and associativity of multiplication by induction.

Use the already declared \(<\) and \(\leq\) abbreviations.  The predecessor
axiom and induction give trichotomy: at a zero argument use that every
nonzero number has a predecessor; at two successor arguments reduce the
comparison to their predecessors.  Addition cancellation and the defining
equations then give preservation of order by addition, and by multiplication
with a positive number.  Cancellation of a positive factor follows: unequal
factors are strictly ordered, and their products remain strictly ordered.
These facts justify subtraction of an ordered pair as the unique \(r\)
in an equation \(a+r=b\).  No subtraction symbol is added to the language.

The least-number principle for any formula follows from induction: induction
on \(b\) proves that if a witness exists at most \(b\), a least such witness
exists at most \(b\).  At \(Sb\), either there was already a witness at
most \(b\), or the new endpoint is least.  Applying this to any supplied
witness proves the unbounded least-number implication.  Bounded maxima follow
by the same induction.  These arguments apply to formulas of arbitrary
complexity; only full \(\PA\) induction is being used.

\begin{annexlemma}[Division inside \(\PA\)]
\label{lem:p05-pa-division}
\(\PA\) proves that, for \(d>0\), every \(a\) has unique \(q,r\) with
\(a=dq+r\) and \(r<d\), and that \(q,r\leq a\).
\end{annexlemma}
\begin{proof}
Induct on \(a\).  At zero take \(q=r=0\).  Given \(a=dq+r\),
trichotomy and \(r<d\) imply either \(Sr<d\) or \(Sr=d\).
For \(Sa\), use respectively \((q,Sr)\) or \((Sq,0)\).
For uniqueness, if \(q<q'\) then \(q'\geq q+1\), whence
\(dq+r<dq+d\leq dq'\leq dq'+r'\), contradicting equality of the
two decompositions.  Reverse the roles for \(q'<q\).  Thus \(q=q'\),
and addition cancellation gives \(r=r'\).  The equation gives \(r\leq a\);
since \(d\geq1\), it also gives \(q\leq dq\leq a\).
\end{proof}

\begin{annexlemma}[The actual pairing and beta graph inside \(\PA\)]
\label{lem:p05-pa-pair-beta}
\(\PA\) proves that its existing pairing graph \(\mathsf P(a,b,w)\)
is total and single-valued in \((a,b)\), has unique inverse components
for every \(w\), and implies \(a,b\leq w\).  It proves
\[
 \forall w\,\forall i\,\exists!v\,\mathsf B(w,i,v),
 \qquad
 \forall w\,\forall i\,\forall v\,
       (\mathsf B(w,i,v)\to v\leq w).
\]
\end{annexlemma}
\begin{proof}
Induction on \(u\) proves unique existence of \(t\) with
\(2t=u(u+1)\), together with \(u\leq t\): use \(t=0\) at zero
and replace \(t\) by \(t+u+1\) at the successor.  Uniqueness uses
cancellation of the positive factor \(2\).  Write \(T(u)\) temporarily
for this unique value; every occurrence can be eliminated by its equation.
The construction proves \(T(u+1)=T(u)+u+1\).  In particular it is strictly
increasing and \(T(u)\geq u\).

For fixed \(a,b\), put \(u=a+b\) and \(w=T(u)+b\).
Then \(2w=u(u+1)+2b\), exactly \(\mathsf P(a,b,w)\).
It has a unique \(w\); moreover \(a\leq u\leq T(u)\leq w\)
and \(b\leq w\).

Conversely, for \(w\) choose the greatest \(u\leq w\) with
\(T(u)\leq w\), using bounded maximization.  Such a \(u\) exists
because \(T(0)=0\).  Since \(T(w+1)>w\), maximality gives
\(T(u)\leq w<T(u+1)\).  Write \(w=T(u)+b\); the last inequality
gives \(b\leq u\).  Write \(u=a+b\).  Then
\(\mathsf P(a,b,w)\).  The intervals
\([T(u),T(u+1))\) are disjoint, so this \(u\), and hence \(b,a\),
is unique.  All uses of \(T\) abbreviate the unique arithmetic relation
already established in this proof.

Decode these \(a,b\) from \(w\), and apply
\cref{lem:p05-pa-division} to \(a\) and \(1+(i+1)b>0\).
The resulting unique remainder \(v\) and quotient satisfy
\(a,b\leq w\), \(q\leq a\), and \(v\leq a\).
These are exactly the bounds in the existing \(\mathsf B\) formula.
Its final marked comparison is equivalent in \(\PA\) to the remainder
bound: if \(v<d\), its witness \(r\) in \(v+Sr=d\) has \(r<d\),
and conversely that equation implies \(v<d\).  Thus this division
description is equivalent to \(\mathsf B(w,i,v)\), including \(b=0\),
when the divisor is \(1\).  Division uniqueness proves both assertions.
\end{proof}

\paragraph{A finite congruence argument with an internal induction.}
For positive \(d\), write \(d\mid a\) for \(\exists q\,(dq=a)\),
and write \(\mathsf{Cong}(r,v,d)\) for
\(\exists q\,\exists q'\,(r+dq=v+dq')\).
Say that two positive numbers are coprime when their only positive common
divisor is \(1\).  These are formula abbreviations used in this proof.
An integer coefficient below is a pair of natural numbers denoting their
difference: equality of two such differences is expressed by cross addition.
Addition and multiplication of pairs are given by
\((a,b)+(c,d)=(a+c,b+d)\) and
\((a,b)(c,d)=(ac+bd,ad+bc)\).  Consequently each signed equation below
is an arithmetic formula with natural variables, justified by the semiring
identities proved above; no integer sort or extra arithmetic axiom is assumed.

In \(\PA\), coprime positive \(A,D\) have signed coefficients
\(s,t\) with \(sA+tD=1\).  To prove this, choose the least positive
signed linear combination \(\delta\) of \(A,D\); the set is nonempty
because it contains \(A\).  Divide \(A\) by \(\delta\).  Its remainder
is another signed linear combination, so minimality makes that remainder
zero.  Dividing \(D\) gives the same conclusion.  Hence \(\delta\)
is a positive common divisor and must equal \(1\).  This proves the
coefficient assertion using the least-number principle and the division
lemma.  It also proves the converse: a common divisor divides the displayed
combination \(1\).

It follows that if \(A,D\) are coprime and \(v<D\), every prescribed
residue \(r\) modulo \(A\) can be combined with residue \(v\) modulo
\(D\).  Indeed reduce the signed number \(s(v-r)\) modulo \(D\) to
\(k<D\), and put \(r'=r+Ak\).  Reduction of a signed coefficient uses
division of its two natural components; choose the nonnegative difference
of their remainders, adding \(D\) if needed.  Then
\(\mathsf{Cong}(r',r,A)\) and \(\mathsf{Cong}(r',v,D)\).
The same signed-coefficient identity shows that if \(A,D\) are each
coprime to \(E\), then \(AD\) is coprime to \(E\): multiply
their two coefficient equations for \(1\).

\begin{annexlemma}[Beta extension in \(\PA\)]
\label{lem:p05-pa-beta-extension}
\(\PA\) proves
\[
\begin{split}
 \forall w\,\forall n\,\forall a\,\exists w'\,\bigl[
 &\mathsf B(w',n,a)\\
 &\land\forall i<n\,\forall v\leq w\,
          (\mathsf B(w,i,v)\to\mathsf B(w',i,v))\bigr].
\end{split}
\]
\end{annexlemma}
\begin{proof}
Work in \(\PA\) with parameters \(w,n,a\).  Induction on \(r\)
proves that a number \(m>w+a\) exists divisible by every positive
\(k\leq r\).  At \(r=0\) take \(w+a+1\); at the successor replace
the preceding positive \(m\) by \((r+1)m\).  Previously required
divisors remain divisors, and the new divisor divides the product.  Apply
this with \(r=n+1\).  This establishes existence before using any
sequence, factorial or product-of-a-sequence function.

For \(i\leq n\), put \(d_i=1+(i+1)m\).  These are pairwise coprime.
For suppose \(i<j\) and a positive \(h>1\) divides \(d_i,d_j\).
Writing \(j=i+q\) gives \(h\mid qm\).  From \(h\mid d_i\)
one has an explicit signed combination of \(h,m\) equal to \(1\),
namely \(d_i-(i+1)m=1\).  Multiplication by \(q\) shows
\(h\mid q\).  Since \(0<q\leq n\), this gives \(h\leq n\),
so the choice of \(m\) gives \(h\mid m\).  It would follow that
\(h\mid d_i-(i+1)m=1\), a contradiction.

For \(i<n\), let the target residue be the unique value supplied by
\(\mathsf B(w,i,v)\); for \(i=n\), let it be \(a\).
Each target is below \(m\), hence below \(d_i\), by
\cref{lem:p05-pa-pair-beta}.  This describes an arithmetic relation with
exactly one target at each \(i\leq n\), not an assumed coded list.

Induct on \(k\leq n+1\) with the following arithmetic invariant:
there exist \(b\) and \(M>0\) such that \(b\) has the target residues
at all \(i<k\), each such \(d_i\) divides \(M\), and \(M\) is
coprime to every \(d_j\) with \(k\leq j\leq n\).
The phrase ``target residues'' expands by the relation in the preceding
paragraph and quantification over its unique value.  Thus this is a single
\(\Lang_A\)-formula with parameters \(w,n,a,m\), licensed for induction.
At zero take \(b=0,M=1\).  At \(k\), use the preceding two-modulus
argument with \(M,d_k\) to obtain
\(b'=b+Mt\) with the required residue at \(k\); set \(M'=Md_k\).
The old congruences persist because every old modulus divides \(M\).
Divisibility by all moduli through \(k\) is immediate.  Coprimality with
the remaining moduli follows from the product property and pairwise
coprimality.  This proves the invariant's successor step.

At \(k=n+1\) obtain \(b\) with all desired residues.  Choose \(w'\)
with \(\mathsf P(b,m,w')\), using \cref{lem:p05-pa-pair-beta}.
Division uniqueness, since every target is below its modulus, identifies
the beta entries of \(w'\) with exactly those targets.  The pairing
component bounds supply every bound in \(\mathsf B\).  This proves the
displayed extension formula.  When \(n=0\), the old residue conditions
are empty and the same argument supplies the one required new entry.
\end{proof}

\begin{annexlemma}[Finite evidence collection in \(\PA\)]
\label{lem:p05-pa-finite-collection}
For every arithmetic formula \(R(i,v,\bar z)\), \(\PA\) proves the
universal closure of
\[
 \forall i<n\,\exists v\,R(i,v,\bar z)
 \quad\longrightarrow\quad
 \exists w\,\forall i<n\,\exists v\leq w\,
      (\mathsf B(w,i,v)\land R(i,v,\bar z)).
\]
\end{annexlemma}
\begin{proof}
Induct on \(n\), with all parameters retained.  At zero take \(w=0\).
At \(Sn\), restrict the premise to \(i<n\), use the induction
hypothesis, and choose a witness \(a\) for \(R(n,a,\bar z)\).
Apply \cref{lem:p05-pa-beta-extension} to append \(a\).
Every old witness remains an entry of the new code; all new entry values
are bounded by that new code by \cref{lem:p05-pa-pair-beta}.
The cases \(i<n\) and \(i=n\) give the conclusion.  Replacing \(n\)
by \(Sm\) gives the version with \(i\leq m\).  No restriction on the
quantifier complexity of \(R\) is needed because the induction scheme is
the full \(\PA\) scheme.
\end{proof}

\paragraph{The graph construction and its internal equations.}
We now prove \cref{thm:p05-pa-pr-graphs} by metainduction on the finite
description \(d\).  The formulas for zero, successor and projections are
atomic and have their stated unique outputs.  In a composition clause,
the induction hypotheses give unique intermediate outputs and the unique
output of \(g\); finite successive existential introductions prove
existence, and uniqueness at each named output proves uniqueness of the
composition.  This argument also covers an empty list of intermediate
outputs.  Its defining equation is precisely its displayed formula.

For the recursion clause, assume the internal totality and uniqueness
theorems for \(\Gamma_g,\Gamma_h\).  First prove existence of an output
for \(\Gamma_d(\bar x,n,y)\) by induction on \(n\) in \(\PA\).
At zero choose the unique base value \(u\) and use beta extension with
\(n=0\) to obtain a code having first entry \(u\).  The step range is
empty, so this code witnesses \(\Gamma_d(\bar x,0,u)\).

At the successor, the induction hypothesis supplies a sequence code \(w\)
through position \(n\) and its last value \(y\).  The step totality
theorem supplies \(z\) with \(\Gamma_h(\bar x,n,y,z)\).
Beta extension at position \(Sn\) preserves positions \(0,\ldots,n\)
and appends \(z\).  All earlier recursion clauses persist, and the new
clause uses the displayed \(\Gamma_h\) fact.  This is an output sequence
for \(\Gamma_d(\bar x,Sn,z)\).  The formula to which induction is
applied is \(\exists y\,\Gamma_d(\bar x,n,y)\).

For uniqueness, take two witnessing codes for the same \(\bar x,n\).
Induct on \(i\leq n\) that their entries at \(i\) agree.  At zero,
both entries satisfy \(\Gamma_g\), so they agree by its uniqueness.
At \(Si\leq n\), the two step clauses have equal input entries by
induction, and uniqueness for \(\Gamma_h\) makes their output entries
equal.  The beta uniqueness lemma identifies these with the actual entries
of the two codes.  At \(i=n\) their outputs agree.  More explicitly,
the induction formula quantifies over the two entry values and says that
the two \(\mathsf B\) facts imply their equality; the two witnessing
codes and \(\bar x,n\) are retained parameters.  This proves the internal
unique-output assertion, rather than uniqueness only for standard histories.

The zero defining equation follows because the first and last entries are
the same entry, with converse furnished by the one-entry construction.
For the successor equation, a code satisfying the length-\(Sn\) condition
also satisfies the length-\(n\) condition when restricted to its earlier
indices.  Its position-\(n\) value and last step give the right-hand side.
Conversely, start with a code for \(\Gamma_d(\bar x,n,y)\) and the
supplied \(\Gamma_h(\bar x,n,y,z)\) fact; the append construction just
proved gives the left-hand side.  No fresh sequence is needed in the
restriction direction, because the graph formula constrains a prefix of
entries rather than imposing a unique numerical code for that prefix.

\paragraph{Syntactic complexity and standard values.}
Every \(\mathsf B\) occurrence expands to the already fixed marked-bounded
formula.  Composition uses existential quantification and conjunction.
Recursion additionally uses the marked bounded universal \(\forall i<n\)
outside existential witnesses.  These are all closure operations of the
declared \(\Sigma^0_1\) class after the fixed negation-normal-form conversion.
Thus every \(\Gamma_d\) is \(\Sigma^0_1\); no exchange of a bounded
universal and an unbounded existential is required.

In \(\mathcal N\), external induction on the description identifies the
initial and composition clauses with their denoted numerical functions.
At recursion, an ordinary induction on \(n\) identifies every stored
entry with the numerical recursion value.  Sequence extension supplies the
converse witnesses.  This proves exactly the stated standard-model graph
equivalence, independently of the internal proof just given.

Finally, for a fixed tuple of standard numerals, evaluate the finite
description using its initial and composition equations and, at every
recursion call, its zero and successor equations a standard finite number
of times.  Induction on the description, with this metainduction on each
standard recursion argument, yields an actual \(\PA\)-derivation of
\(\Gamma_d(\overline{\bar a},\overline{f_d(\bar a)})\).
The proved universal uniqueness theorem then gives the displayed
numeral-by-numeral functional representation.  This last conclusion is in
\(\PA\); it does not replace the separately cited \(\mathsf Q\)
representation theorem.

The whole argument has used a separately fixed finite description and a
finite metainduction on it.  It has not introduced a uniformly total
evaluator for all descriptions, nor has it verified the particular
proof-constructor invariants needed for internal modus ponens and positive
introspection.  Those applications require their own specified descriptions
and arithmetic proofs.

\paragraph{Bounded evidence, including dependent witnesses.}
We prove \cref{thm:p05-pa-sigma-evidence}.  First take the fixed
negation-normal-form conversion of \(\varphi\).  Each propositional or
quantifier-duality step is a classical first-order equivalence, and its
marked bounded version follows by expanding the declared bounded
quantifier.  Thus \(\PA\) proves equivalence with this conversion.
For the resulting \(\Sigma^0_1\) formula define evidence by the following
recursion, taking the first clause whenever the whole subformula is bounded.
This priority removes the possible overlap between the first clause and
the connective clauses.  In the displays, \(\bar x\) comprises the free
parameters other than the bound variable explicitly shown.
Choose the evidence variable fresh from the original formula; all subsequent
auxiliary bound variables are freshened before the displayed substitutions.
\begin{enumerate}
\item For bounded \(\theta\), put
\(\mathsf{Ev}_\theta(\bar x,e)\defeq e=0\land\theta(\bar x)\).
\item For \(\psi\land\chi\), put
\[
 \mathsf{Ev}_{\psi\land\chi}(\bar x,e)\defeq
 \exists u\leq e\,\exists v\leq e\,
 (\mathsf P(u,v,e)\land\mathsf{Ev}_\psi(\bar x,u)
                         \land\mathsf{Ev}_\chi(\bar x,v)).
\]
For \(\psi\lor\chi\), put
\[
\begin{split}
 \mathsf{Ev}_{\psi\lor\chi}(\bar x,e)\defeq
 \exists d<\bar2\,\exists v\leq e\,\bigl[
 &\mathsf P(d,v,e)\\
 &\land((d=0\land\mathsf{Ev}_\psi(\bar x,v))
       \lor(d=\bar1\land\mathsf{Ev}_\chi(\bar x,v)))\bigr].
\end{split}
\]
\item For an unbounded existential, put
\[
 \mathsf{Ev}_{\exists y\,\psi}(\bar x,e)\defeq
 \exists y\leq e\,\exists v\leq e\,
   (\mathsf P(y,v,e)\land\mathsf{Ev}_\psi(y,\bar x,v)).
\]
For \(\exists y<t(\bar x)\,\psi\), replace the first displayed
bound by \(y<t(\bar x)\).  For a \(\leq\)-marked existential use
that marked bound instead.  Pairing itself supplies \(y\leq e\)
whenever the displayed matrix holds.
\item For a marked bounded universal, put
\[
 \mathsf{Ev}_{\forall y<t(\bar x)\,\psi}(\bar x,e)\defeq
 \forall y<t(\bar x)\,\exists v\leq e\,
   (\mathsf B(e,y,v)\land\mathsf{Ev}_\psi(y,\bar x,v)).
\]
Use the corresponding \(\leq\)-marked bound for that variant.
\end{enumerate}
Every quantifier introduced in these clauses is marked bounded, every bound
is an arithmetic term independent of its bound variable, and
\(\mathsf P,\mathsf B\) have their fixed quantifier-free or bounded
expansions.  Induction on formation proves that each evidence formula is
bounded.  It does not require moving an unbounded quantifier across a bound.

Prove in \(\PA\), by metainduction on the fixed formation, both
\(\mathsf{Ev}_\psi(\bar x,e)\to\psi(\bar x)\) and
\(\psi(\bar x)\to\exists e\,\mathsf{Ev}_\psi(\bar x,e)\).
In the bounded case the first implication eliminates a conjunction, and
the second uses \(e=0\).  For a conjunction, forward elimination and
the two induction hypotheses give its conjuncts.  Conversely choose the
two evidence values supplied by induction and pair them; internal pairing
totality and its component bounds supply the required bounded witnesses.
For a disjunction, the evidence's two tagged alternatives give the relevant
disjunct by induction.  Conversely use tag \(0\) or \(1\) with the
evidence for the given disjunct, and pair the tag with that evidence.

For an existential, forward elimination supplies its witness and evidence
for its matrix; apply induction and then existential introduction.
Conversely choose the original witness, obtain its matrix evidence by
induction, and pair the two values.  The pairing bounds give the witnesses
in the unbounded-existential evidence clause.  In the marked bounded case,
the original witness already satisfies the required marked bound.

For a marked bounded universal, forward elimination supplies evidence for
each index in the bound; induction yields the matrix at that index and
universal introduction gives the original formula.  Conversely apply the
induction hypothesis at every bounded index to obtain one existential
evidence value for that index.  \Cref{lem:p05-pa-finite-collection}, with
\(R(y,v,\bar x)=\mathsf{Ev}_\psi(y,\bar x,v)\) and
\(n=t(\bar x)\), gives one beta code \(e\) with exactly the required
property.  Use its \(\leq\) variant when needed.  At an empty strict
bound, \(e=0\) suffices, consistently with that lemma's base case.
The bound variable is absent from its bounding term, so substitution of
the term for \(n\) is legitimate throughout.

These are internal derivations for the fixed formulas with all displayed
parameters.  They establish the theorem for the converted formula; composing
with its first-order equivalence to \(\varphi\) establishes the claimed
result for \(\varphi\). Bounded evidence and an ND proof code are
different objects. The separate construction in
\cref{thm:p05-original-sigma-evidence} turns this evidence into an annotated
derivation and proves its selected checker invariant inside \(\PA\).

\begin{annexlemma}[An exact numerical modus-ponens constructor]
\label{lem:p05-numerical-mp-constructor}
There is a specified primitive recursive function \(\mathsf{mp}(p,q)\),
independent of \(T\), such that for every fixed c.e. axiom set \(T\)
and sentences \(\varphi,\psi\), the numerical relations satisfy
\[
\begin{split}
 &\mathsf{Proof}_T(p,\ulcorner\varphi\to\psi\urcorner)
       \land\mathsf{Proof}_T(q,\ulcorner\varphi\urcorner)\\
 &\hspace{12mm}\Longrightarrow
       \mathsf{Proof}_T(\mathsf{mp}(p,q),\ulcorner\psi\urcorner).
\end{split}
\]
This statement is about the certified numerical relation of
\cref{sec:arithmetic-coding-04}; it is not yet an internal \(\PA\) formula.
\end{annexlemma}
\begin{proof}
Here is the constructor, with the argument order fixed by the display.
Decode \(p=N_{16}(N_{15}(R_p),W_p)\) and
\(q=N_{16}(N_{15}(R_q),W_q)\).  Check both derivation records and their
sentence roots, and that each witness list has the length of its root open
list; do not perform any axiom-membership test.  Check that the root of
\(R_p\) is an implication whose antecedent is the root of \(R_q\).
Return zero if any of these checks fails.

Otherwise let \(a=\ell(R_q)\), \(b=\ell(R_p)\), and let \(z\)
be the consequent code.  In each row \([r,u,J,U,l]\) of \(R_p\),
replace each entry \(j\) of the premise list \(J\) by \(a+j\),
leaving the other four fields unchanged.  Call the resulting list
\(R_p^{+a}\).  Set
\[
\begin{aligned}
 R&=R_q\mathbin{+\!+}R_p^{+a}\mathbin{+\!+}
          [[8,z,[a-1,a+b-1],0,0]],\\
 W&=W_q\mathbin{+\!+}W_p,\\
 \mathsf{mp}(p,q)&=N_{16}(N_{15}(R),W),
\end{aligned}
\]
where \(+\!+\) denotes list concatenation only in this display.
Both lengths are positive because the derivation checks succeeded, so the
two decremented indices exist.  The tag \(8\) is the actual body-order
implication-elimination rule: antecedent first, implication second.
The tests, row shift and concatenations are the primitive recursive
operations already constructed for these codes; every loop is bounded by a
decoded list length.  Thus this prescription specifies a total primitive
recursive function, including its failure value.

For valid input proofs, all old parent counts and earlier-reference conditions
are preserved inside each copied tree.  Each old root receives exactly one
new parent, the appended row.  Its premise formulas satisfy the actual
implication-elimination rule, and it discharges nothing.  All parameters
in \(U\) are formula/term codes, variable indices or discharge labels,
not row references, so they require no shift.  Assumption markers and their
matching discharge labels also remain unchanged: labels in different input
trees may coincide, but the only new common ancestor is a nondischarge rule.
No occurrence in the other tree lies above an old discharge application.

The new open list is the old list for \(R_q\), followed by the shifted
old list for \(R_p\).  Consequently \(W_q\) followed by \(W_p\)
places exactly the old accepting certificate at each surviving leaf.
Every \(V(e_T,\cdot,\cdot)\) call is therefore an unchanged successful
input call.  The consequent is a sentence because both subformulas of the
sentence \(\varphi\to\psi\) are sentences.  All numerical proof checks
now succeed, including the case of empty open and witness lists.
\end{proof}
\end{annexitem}

\begin{proofboundary}{def:p05-internal-proof-predicate-boundary}{The selected standard PA proof construction}
The preceding library supplies the arithmetic and bounded evidence
collection used by the subsequent direct proof construction. Its selected
formulas are
\[
 \rho(p,x)\defeq\mathsf{Prf}^{19}_{\PA}(p,x),\qquad
 P_\rho(x)\defeq\exists p\,\rho(p,x)
                       =\mathsf{Prv}^{19}_{\PA}(x).
\]
\Cref{def:p05-standard-pa-numerical-proof,def:p05-selected-direct-pa-graph}
fixes the direct axiom convention, the actual twenty-rule checker,
its expanded finite description and the literal \(\Gamma\) formula.
\Cref{thm:p05-standard-pa-numerical-proof,thm:p05-selected-direct-pa-graph}
proves its standard-model proof meaning. The certified c.e. code convention
above remains a separate presentation.

The direct constructor \(\mathsf{mp}_{19}\) in
\cref{def:p05-direct-mp-description} has the internal invariant proved in
\cref{thm:p05-direct-pa-internal-mp}. For every fixed pair of sentences
\(\varphi,\psi\), it gives in \(\PA\) the universal closure in
\(p,q\) of
\[
\begin{split}
 &\rho(p,\overline{\ulcorner\varphi\to\psi\urcorner})
       \land\rho(q,\overline{\ulcorner\varphi\urcorner})\\
 &\hspace{8mm}\to
       \rho(\mathsf{mp}_{19}(p,q),\overline{\ulcorner\psi\urcorner}),
\end{split}
\]
where the function application is eliminated through its specified
\(\Gamma_{d_{\rm mp}}\) graph. Its internal row, parent-count and
open-leaf invariants concern the direct tag-19 records. The numerical
certified-code \(\mathsf{mp}\) above is a different constructor.
Existential elimination and introduction give D2 for \(P_\rho\), and
the same theorem supplies its D1 metatheoretic implication.

For every fixed closed \(\Sigma^0_1\) sentence \(\sigma\) in its
original syntax, \cref{thm:p05-original-sigma-evidence} supplies the
specified unary description \(\widehat d_\sigma^+\). Writing
\(B_\sigma\) for its numerical output function, the theorem proves
in \(\PA\)
\[
 \forall e\,(\mathsf{Ev}_\sigma(e)\to
       \rho(B_\sigma(e),\overline{\ulcorner\sigma\urcorner})),
\]
with the function application eliminated through
\(\Gamma_{\widehat d_\sigma^+}\). Its signed construction gives the
original formula code even when the evidence was defined through marked
negation normal form. Combining it with \cref{thm:p05-pa-sigma-evidence}
gives
\(\PA\derives\sigma\to P_\rho(\overline{\ulcorner\sigma\urcorner})\).
Taking \(\sigma=P_\rho(\overline{\ulcorner\varphi\urcorner})\)
gives D3, with the literal graph and outer sentence-code construction
verified in \cref{thm:p05-literal-graph-evidence,thm:p05-direct-pa-d3}.
These later proofs use the library established above; the library proofs
do not assume D3. The result concerns this selected arithmetic formula
and makes no internal identification with an arbitrary numeralwise
representative. Further provability-logic completeness results require
their separately stated constructions.
\end{proofboundary}
